\documentclass[12pt,a4paper,twoside]{article}

% ─── Packages ────────────────────────────────────────────────────────────────
\usepackage[utf8]{inputenc}
\usepackage[T1]{fontenc}
\usepackage[english]{babel}
\usepackage{geometry}
\usepackage{graphicx}
\usepackage{amsmath}
\usepackage{amssymb}
\usepackage{booktabs}
\usepackage{multirow}
\usepackage{caption}
\usepackage{subcaption}
\usepackage{hyperref}
\usepackage{listings}
\usepackage{xcolor}
\usepackage{float}
\usepackage{tikz}
\usepackage{pgf-umlsd}
\usepackage{setspace}
\usepackage{fancyhdr}
\usepackage{natbib}
\usepackage{algorithm}
\usepackage{algpseudocode}
\usepackage{url}

\geometry{
  a4paper,
  top=2.5cm, bottom=2.5cm,
  left=3cm, right=2.5cm
}

\onehalfspacing

\hypersetup{
  colorlinks=true,
  linkcolor=black,
  citecolor=black,
  urlcolor=blue
}

% ─── Title info ──────────────────────────────────────────────────────────────
\title{
  \vspace{-1cm}
  \Large\textbf{AI-Based Decision Support for Goalkeeper\\
  Goal-Line Infringements in Football Penalty Kicks}\\[0.5cm]
  \large Bachelor Thesis
}
\author{[Your Name]}
\date{2026}

% ─── Document ────────────────────────────────────────────────────────────────
\begin{document}

\maketitle
\thispagestyle{empty}

\begin{abstract}
Goalkeeper goal-line infringements during penalty kicks are a specific class of
football rule violations subject to Video Assistant Referee (VAR) review.
Determining whether the goalkeeper's foot crossed the goal line at the moment of
the kick requires precise spatial and temporal analysis of broadcast footage,
which is currently performed manually by officials.

This thesis presents a decision-support pipeline that automates this analysis
from single-camera broadcast video.
The system performs automatic kick-moment detection via ball motion analysis,
extracts the relevant video frame, applies a fine-tuned YOLOv8n detector to
localize the goalkeeper and the ball, reconstructs the goal-line geometry using
Hough-based line detection, and computes the goalkeeper's foot-to-line distance.
An explicit uncertainty policy converts geometrically ambiguous cases into an
\textit{uncertain} output, enabling the system to abstain rather than commit to
a potentially wrong decision.

Evaluated on 22 held-out test clips, the final pipeline achieves a coverage of
0.909, a selective accuracy of 0.900, and a violation recall of 1.000 with a
precision of 0.800, yielding an F1 score of 0.889 for the violation class.
No confirmed violation is missed among the clips for which the system
produces a definite decision.

Additional experiments compare a YOLOv8n baseline against a larger YOLO26n
detector and investigate pose-estimation-based foot localization as a refinement
strategy. Both alternatives improved component-level metrics but did not improve
the end-to-end pipeline result, illustrating that detector quality alone does not
determine decision quality. An experimental extension to player encroachment
detection is also presented as a demonstration of the pipeline's generalisability.

The thesis contributes a complete end-to-end methodology, an annotated dataset
split, and an evaluation framework for broadcast-based penalty-kick officiating
support.
\end{abstract}

\newpage
\tableofcontents
\newpage

% ═════════════════════════════════════════════════════════════════════════════
\section{Introduction}
\label{sec:intro}
% ═════════════════════════════════════════════════════════════════════════════

\subsection{Motivation and Problem Statement}

Football penalty kicks are among the highest-stakes moments in the sport.
They determine outcomes of knockout rounds, championships, and relegation battles.
At the same time, the rules governing penalty kicks are precise: the goalkeeper
must remain on the goal line, with at least one foot touching or behind the line,
until the ball is kicked~\cite{ifab2024}.
A goalkeeper who steps forward before the kick gains a positional advantage that
is considered unfair under the Laws of the Game.

Since 2018, Video Assistant Referee (VAR) technology has been used at the
highest levels of football to review potential rule violations, including
goalkeeper encroachment during penalty kicks.
However, VAR decisions on goalkeeper position rely on human operators reviewing
video manually. This process is time-consuming, subjective, and inconsistent:
the threshold for what counts as a ``clear and obvious error'' varies across
matches and competitions.

Automating the detection of goalkeeper goal-line infringements from broadcast
video would provide a consistent, objective, and rapid decision-support signal
to officiating teams. This thesis addresses that problem.

\subsection{Scope and Constraints}

The problem is constrained to broadcast-quality video footage:
\begin{itemize}
  \item Single fixed camera, standard broadcast angle
  \item 720p resolution at approximately 25 frames per second
  \item No depth sensors, calibrated cameras, or tracking data
  \item No assumption of camera calibration or fixed pitch geometry
\end{itemize}

These constraints reflect the realistic deployment scenario, where broadcast
footage is the only available data source for post-hoc VAR review.

\subsection{Approach Overview}

The proposed system operates as a sequential pipeline:
\begin{enumerate}
  \item A penalty kick clip is provided as input.
  \item The kick moment is detected automatically from ball motion.
  \item The relevant frame (one frame before the detected kick) is extracted.
  \item A fine-tuned YOLO detector localises the goalkeeper and ball.
  \item The goal line is detected geometrically from the extracted frame.
  \item The goalkeeper's foot position relative to the line is evaluated.
  \item An uncertainty policy decides whether to commit to a decision or abstain.
\end{enumerate}

The system is framed as a \textbf{decision-support tool}, not as a replacement
for human officiating. The explicit uncertainty output is a design choice: in
cases where the geometry is too ambiguous for a reliable decision, the system
signals that human review is required rather than producing a potentially
wrong answer.

\subsection{Contributions}

The main contributions of this thesis are:

\begin{itemize}
  \item A complete end-to-end pipeline for goalkeeper goal-line infringement
        detection from broadcast video.
  \item An annotated dataset of 135 penalty kick clips with violation / valid /
        uncertain labels and a train/validation/test split.
  \item A quantitative evaluation framework including abstaining classifier
        metrics appropriate for decision-support tools.
  \item An analysis of the interplay between detector quality and end-to-end
        decision quality, showing that component improvement does not guarantee
        pipeline improvement.
  \item An experimental extension to player encroachment detection, demonstrating
        the generalisability of the pipeline architecture.
\end{itemize}

\subsection{Thesis Structure}

Section~\ref{sec:related} reviews related work. Section~\ref{sec:dataset}
describes the dataset and annotation process. Section~\ref{sec:method} presents
the methodology in detail. Section~\ref{sec:eval} defines the evaluation
framework. Section~\ref{sec:results} reports the experimental results.
Section~\ref{sec:discussion} discusses the findings and limitations.
Section~\ref{sec:conclusion} concludes and outlines future work.


% ═════════════════════════════════════════════════════════════════════════════
\section{Related Work}
\label{sec:related}
% ═════════════════════════════════════════════════════════════════════════════

\subsection{Sports Video Analysis}

Automated analysis of sports footage has been an active research area for two
decades. Early work focused on scene segmentation and event
detection~\cite{ekin2003automatic}.
Convolutional neural networks enabled more reliable player and ball
tracking~\cite{suda2021football}, and later transformer-based architectures
improved multi-object tracking in crowded sports scenes~\cite{yu2022sportsmot}.

Ball detection is a well-studied sub-problem with specific
challenges~\cite{theagarajan2018soccer}: the ball is small, fast, frequently
occluded, and subject to motion blur. The TrackNetV2 architecture was
specifically designed for high-speed ball tracking from video~\cite{huang2019tracknet}.
YOLOv8-based approaches have also been applied to real-time ball and player
detection in broadcast footage~\cite{ultralytics2023}.

\subsection{Rule Violation Detection}

Goal-line technology (GLT) is the established industrial solution for
detecting whether the ball has fully crossed the goal line. Systems such as
Hawk-Eye and GoalControl use multi-camera setups with calibrated cameras and
real-time tracking~\cite{hawk2012}. These are dedicated hardware installations
incompatible with post-hoc broadcast analysis.

For goalkeeper encroachment specifically, there is limited prior academic work.
Sporadic references appear in surveys of VAR-related decision support, but no
public dataset or pipeline has been previously published for this specific task.
This thesis represents an initial investigation into the problem under realistic
broadcast constraints.

\subsection{Pose Estimation in Sports}

Human pose estimation has been applied to sports biomechanics and action
recognition~\cite{bridgeman2019multi}. YOLOv8-pose offers real-time
multi-person pose estimation from single images~\cite{ultralytics2023}.
The use of pose keypoints for foot localization is a natural extension of
goalkeeper analysis, and it is investigated in this thesis as an alternative
to bounding-box-based foot proxies. MediaPipe Pose provides an alternative
lightweight pose estimation
option~\cite{lugaresi2019mediapipe}.

\subsection{Uncertainty in Classification}

Abstaining classifiers, also called \textit{selective classifiers} or
\textit{classifiers with reject option}, explicitly model the option of
withholding a prediction when confidence is low~\cite{bartlett2008classification}.
This framework is particularly relevant for high-stakes decision support,
where the cost of a wrong confident prediction exceeds the cost of abstaining.
The evaluation metrics used in this thesis (coverage, selective accuracy,
lower-bound accuracy) follow the abstaining classifier
literature~\cite{geifman2017selective}.


% ═════════════════════════════════════════════════════════════════════════════
\section{Dataset}
\label{sec:dataset}
% ═════════════════════════════════════════════════════════════════════════════

\subsection{Data Collection}

The dataset was compiled from publicly available match broadcast footage of
professional football matches. Penalty kick moments were identified manually
by reviewing match recordings. For each penalty kick, a short clip of
approximately 10 seconds centred around the kick moment was extracted at 720p
resolution and 25~fps.

In total, 135 penalty kick clips were collected from matches spanning multiple
European leagues and competitions (2014--2017). The clips vary in:
\begin{itemize}
  \item Broadcast angle and zoom level
  \item Pitch conditions and lighting
  \item Goalkeeper stance and movement
  \item Degree of goalkeeper goal-line compliance
\end{itemize}

\subsection{Annotation}

Each clip was annotated with a ternary label:
\begin{itemize}
  \item \textbf{violation}: the goalkeeper's foot was clearly forward of the
        goal line at the moment of the kick.
  \item \textbf{valid}: the goalkeeper's foot was on or behind the goal line
        at the moment of the kick.
  \item \textbf{uncertain}: the camera angle, occlusion, or motion blur made
        it impossible to determine the goalkeeper's position reliably.
\end{itemize}

Annotation was performed by manual frame-by-frame review. The kick moment
was defined as the frame at which the ball begins to move, identified by visual
inspection and confirmed by ball motion analysis.

Ambiguous cases were labeled \textit{uncertain} by consensus. These clips are
excluded from binary precision/recall evaluation but are retained in the dataset
for completeness.

\subsection{Train / Validation / Test Split}

The 135 clips were split into three non-overlapping sets with no match overlap
across splits (i.e., no two clips from the same match appear in different splits):

\begin{table}[H]
  \centering
  \caption{Dataset split statistics}
  \label{tab:splits}
  \begin{tabular}{lrrr}
    \toprule
    Split & Clips & Violations & Valid \\
    \midrule
    Train & 94 & -- & -- \\
    Validation & 19 & -- & -- \\
    Test & 22 & 4 & 16 (+2 uncertain) \\
    \midrule
    Total & 135 & & \\
    \bottomrule
  \end{tabular}
\end{table}

The test split of 22 clips contains 4 confirmed violations, 16 confirmed valid
cases, and 2 clips labeled \textit{uncertain} (excluded from violation-focused
evaluation). All quantitative results reported in this thesis use the test split
exclusively; the train and validation splits were used for model training and
hyperparameter selection only.

\subsection{YOLO Detection Dataset}

For training the YOLO object detector, a subset of frames from the penalty kick
clips was annotated with bounding boxes for two classes: goalkeeper and ball.
The resulting detection dataset contains:
\begin{itemize}
  \item Training images: 354
  \item Validation images: 66
  \item Test images: 88
\end{itemize}

Each frame was extracted at multiple temporal offsets from the kick moment to
ensure coverage of different goalkeeper positions. Bounding boxes were drawn
manually around the goalkeeper (full body) and ball (when visible) in each frame.


% ═════════════════════════════════════════════════════════════════════════════
\section{Methodology}
\label{sec:method}
% ═════════════════════════════════════════════════════════════════════════════

\subsection{Pipeline Overview}

The full system processes a penalty kick video clip through the following stages:

\begin{enumerate}
  \item \textbf{Automatic kick detection}: ball motion is analysed to identify
        the kick frame.
  \item \textbf{Frame adjustment}: the detected frame is shifted by $-1$ to
        correct for systematic late-detection bias.
  \item \textbf{Frame extraction}: the adjusted kick frame is saved as an image.
  \item \textbf{YOLO detection}: a fine-tuned YOLOv8n model detects the
        goalkeeper and ball in the extracted frame.
  \item \textbf{Goal-line geometry}: a Hough-based line detector reconstructs
        the goal line from the image.
  \item \textbf{Distance computation}: the goalkeeper's foot proxy is projected
        onto the goal line and the signed distance is computed.
  \item \textbf{Uncertainty policy}: a rule-based policy converts geometrically
        ambiguous decisions to \textit{uncertain}.
  \item \textbf{Decision output}: the final decision (on\_line / off\_line /
        uncertain) is produced with supporting metadata.
\end{enumerate}

\subsection{Automatic Kick Detection}

\subsubsection{Ball Tracking}

The kick moment is estimated from ball motion within a temporal window of
$[0.5\text{s},\ 2.5\text{s}]$ from the start of the input clip. This window
was chosen to capture the kick in a standard 10-second penalty clip with some
margin for kick timing variability.

For each frame in the window, the YOLOv8n detector is applied and the highest-
confidence ball detection is retained (minimum confidence threshold: 0.3).
Frames without a valid ball detection, or where the ball position jumps by more
than 180 pixels between consecutive frames (tracking jump filter), are excluded.

\subsubsection{Velocity Analysis and Onset Detection}

From the retained ball positions, a frame-to-frame displacement sequence is
computed. The baseline velocity is estimated as the median velocity over a
stable pre-kick period. The kick moment is detected as the first frame where
the ball velocity exceeds a threshold of $2.5 \times$ the baseline velocity for
a sustained period (velocity prominence threshold).

If no onset is found, the system falls back to the frame of peak ball velocity
(peak fallback). If ball detections are insufficient for reliable velocity
estimation (fewer than 5 retained frames), kick detection fails and the clip
is marked as a pipeline failure.

\subsubsection{Frame Adjustment}

Analysis of detection errors on the validation set revealed a systematic
late-detection bias: the detected kick frame was on average 2.1 frames later
than the annotated kick frame, with a median error of 1 frame. To correct for
this bias, a fixed $-1$ frame adjustment is applied to the detected kick frame
before further processing.

This adjustment is the only temporal correction applied. It was selected after
ablation over offsets $\{-2, -1, 0, +1\}$ by maximising the end-to-end pipeline
exact match rate on the validation set.

\subsection{Goalkeeper and Ball Detection}

A YOLOv8n model (\texttt{train4}) was fine-tuned on the annotated detection
dataset. The model was trained with:
\begin{itemize}
  \item Input resolution: $640 \times 640$ pixels
  \item 100 training epochs
  \item AdamW optimiser with cosine learning rate schedule
  \item Data augmentation: random horizontal flip, mosaic, scale jitter
\end{itemize}

Two detection classes were used: goalkeeper (class 0) and ball (class 1).
The goal-line decision uses only goalkeeper detections; ball detections are
used only by the kick detection module.

For the pipeline, the goalkeeper bounding box with the highest confidence score
is selected. If no goalkeeper detection exceeds the confidence threshold of 0.05,
the clip is assigned an \textit{uncertain} decision with reason \texttt{no\_goalkeeper}.

\subsection{Goal-Line Geometry}

\subsubsection{Line Detection}

The goal line is detected using a pipeline of:
\begin{enumerate}
  \item Restrict analysis to the lower third of the frame (where the goal line
        is expected to appear).
  \item Apply white-pixel masking: pixels with HSV saturation below 30 and
        value above 200 are retained.
  \item Apply Canny edge detection to the masked image.
  \item Run the probabilistic Hough line transform to detect line segments.
  \item Filter: retain only near-horizontal segments ($|\theta| < 15°$) with
        length above 50 pixels.
  \item Fit a single representative line through the retained segments using
        weighted least-squares regression.
\end{enumerate}

If no line segments pass the filter, the clip is assigned \textit{uncertain}
with reason \texttt{no\_line}.

\subsubsection{Foot Proxy Extraction}

The goalkeeper's foot position is approximated by the bottom-centre point of
the goalkeeper bounding box. This is a deliberate simplification: the bottom
of the bounding box corresponds to the lowest visible part of the goalkeeper's
body, which in most upright or semi-crouched postures is a foot or lower shin.
The limitations of this proxy are discussed in Section~\ref{sec:discussion}.

\subsubsection{Distance Computation}

Let $\mathbf{p} = (x_p, y_p)$ be the foot proxy point and let the goal line be
described by $ax + by + c = 0$ (unit normal). The signed pixel distance is:

\begin{equation}
  d = \frac{a x_p + b y_p + c}{\sqrt{a^2 + b^2}}
\end{equation}

A positive distance indicates the goalkeeper's foot is in front of the goal
line (i.e., a potential violation). The raw decision is:
\begin{align*}
  d > 0 &\Rightarrow \texttt{off\_line} \quad \text{(violation)} \\
  d \leq 0 &\Rightarrow \texttt{on\_line} \quad \text{(valid)}
\end{align*}

An additional quantity, \texttt{local\_y\_err\_px}, measures the residual
vertical error of the Hough fit at the projected foot location. High values
indicate that the line-fitting was imprecise at that position.

\subsection{Uncertainty Policy}

The raw binary decision is post-processed by a rule-based uncertainty policy.
The policy converts a decision to \textit{uncertain} when either of two
conditions holds:

\begin{description}
  \item[\texttt{near\_boundary}:] $|d| < \delta_{\text{margin}}$,
        where $\delta_{\text{margin}} = 2.0$~px. At this distance,
        sub-pixel detection noise could flip the decision.
  \item[\texttt{high\_local\_y\_err}:] $e_y > \delta_{\text{yerr}}$,
        where $\delta_{\text{yerr}} = 8.0$~px. This indicates that the
        goal-line geometry is too imprecise at the foot location to support
        a confident decision.
\end{description}

When both conditions hold simultaneously, the reason is reported as
\texttt{boundary\_plus\_local\_y}.

The thresholds $\delta_{\text{margin}}$ and $\delta_{\text{yerr}}$ were
set based on analysis of the detection noise profile on the training split
and were not tuned on test labels.


% ═════════════════════════════════════════════════════════════════════════════
\section{Evaluation Framework}
\label{sec:eval}
% ═════════════════════════════════════════════════════════════════════════════

\subsection{Abstaining Classifier Metrics}

Standard binary classification metrics (accuracy, precision, recall) are
not directly applicable to a system that can abstain on uncertain cases.
We use the abstaining classifier framework~\cite{geifman2017selective} with
the following metrics:

\begin{description}
  \item[\textbf{Coverage}] The fraction of clips for which the system produces
        a definite (non-uncertain) decision:
        $\text{coverage} = n_{\text{certain}} / n_{\text{total}}$.

  \item[\textbf{Selective accuracy}] Accuracy computed only on the clips for
        which the system committed to a decision:
        $\text{sel\_acc} = n_{\text{correct certain}} / n_{\text{certain}}$.

  \item[\textbf{Lower-bound accuracy}] Worst-case accuracy if all abstentions
        are counted as errors:
        $\text{lb\_acc} = n_{\text{correct certain}} / n_{\text{total}}$.

  \item[\textbf{Precision} (violation class)] Among clips where the system
        predicted \textit{violation}, the fraction that are true violations.

  \item[\textbf{Recall} (violation class)] Among all true violations, the
        fraction that were predicted as \textit{violation}.

  \item[\textbf{F1} (violation class)] Harmonic mean of violation precision
        and recall.
\end{description}

\subsection{No-Skill Baselines}

Two degenerate baselines are used for comparison:
\begin{itemize}
  \item \textbf{Always valid}: always predicts \textit{on\_line}.
        Recall (violation) = 0.000.
  \item \textbf{Always violation}: always predicts \textit{off\_line}.
        Recall (violation) = 1.000, precision = class prior = 4/20 = 0.200.
\end{itemize}

A useful system must strictly dominate both baselines.

\subsection{Kick-Timing Metrics}

For the kick detection submodule, standard timing accuracy metrics are used:
\begin{itemize}
  \item Success rate: fraction of clips with a valid kick frame estimate.
  \item Exact accuracy: fraction with detected frame equal to annotated frame.
  \item Within $\pm k$ frames: fraction with $|$detected $-$ annotated$| \leq k$.
  \item Mean/median absolute error (MAE/MedAE) in frames.
\end{itemize}


% ═════════════════════════════════════════════════════════════════════════════
\section{Experimental Results}
\label{sec:results}
% ═════════════════════════════════════════════════════════════════════════════

\subsection{YOLO Detection Results}

Table~\ref{tab:yolo_detection} shows the detection performance of the
fine-tuned YOLOv8n model (\texttt{train4}) on the 88-image test split.

\begin{table}[H]
  \centering
  \caption{YOLOv8n (train4) detection metrics on test split (88 images)}
  \label{tab:yolo_detection}
  \begin{tabular}{lrrrr}
    \toprule
    Class & Precision & Recall & mAP50 & mAP50-95 \\
    \midrule
    Goalkeeper & 0.992 & 1.000 & 0.995 & 0.685 \\
    Ball       & 0.826 & 0.558 & 0.631 & 0.240 \\
    \midrule
    \textbf{Overall} & \textbf{0.909} & \textbf{0.779} & \textbf{0.813} & \textbf{0.462} \\
    \bottomrule
  \end{tabular}
\end{table}

Goalkeeper detection is near-perfect: mAP50 of 0.995 with recall of 1.000
on the test split means no goalkeeper was missed in any evaluated frame.
Ball detection is considerably weaker at mAP50 of 0.631, reflecting the
difficulty of detecting a small, fast-moving object in broadcast footage.
The lower ball recall (0.558) is the primary driver of occasional kick
detection failures, as ball positions are required for velocity analysis.

Since the goal-line decision depends on goalkeeper localisation rather than
ball localisation, the very high goalkeeper detection quality is the critical
result for the main pipeline. Ball detection quality primarily affects the
kick timing submodule.

\subsection{Automatic Kick Detection Results}

Table~\ref{tab:kick_detection} reports kick detection performance evaluated
on 91 penalty clips with annotated kick times.

\begin{table}[H]
  \centering
  \caption{Kick detection metrics (train4 detector, frame\_adjust = $-1$, 91 clips)}
  \label{tab:kick_detection}
  \begin{tabular}{lr}
    \toprule
    Metric & Value \\
    \midrule
    Clips attempted & 91 \\
    Successful detections & 84 (92.3\%) \\
    Exact frame accuracy & 0.264 \\
    Within $\pm 1$ frame & 0.582 \\
    Within $\pm 2$ frames & 0.725 \\
    Within $\pm 3$ frames & 0.758 \\
    Mean absolute error & 5.01 frames \\
    Median absolute error & 1.0 frame \\
    Mean signed error & $-2.1$ frames \\
    \bottomrule
  \end{tabular}
\end{table}

The module successfully estimates the kick frame for 92.3\% of clips.
Among successful detections, the median absolute error is exactly 1 frame
(40\,ms at 25\,fps). The mean signed error of $-2.1$ frames indicates that
after the $-1$ correction, the remaining errors are mixed, with slightly more
early predictions (40) than late predictions (20). The mean absolute error
of 5.01 frames reflects a small number of outlier clips where the velocity
onset was difficult to detect; these outliers inflate the mean without
affecting the median.

The dominant detection method was \texttt{motion\_onset} (73 clips), with
\texttt{velocity\_peak\_fallback} used for 11 clips where no clear onset was
found. 7 clips failed due to insufficient ball detections.

\subsection{End-to-End Pipeline Results}

Table~\ref{tab:final_pipeline} reports the end-to-end evaluation of the final
adopted pipeline on the 22-clip test split.

\begin{table}[H]
  \centering
  \caption{Final pipeline evaluation on test split (22 clips)}
  \label{tab:final_pipeline}
  \begin{tabular}{lr}
    \toprule
    Metric & Value \\
    \midrule
    Clips attempted & 22 \\
    Pipeline successful & 21 \\
    Certain predictions & 20 \\
    Uncertain predictions & 2 \\
    \midrule
    Coverage & 0.909 \\
    Selective accuracy & 0.900 \\
    Lower-bound accuracy & 0.818 \\
    \midrule
    Precision (violation) & 0.800 \\
    Recall (violation) & \textbf{1.000} \\
    F1 (violation) & 0.889 \\
    \midrule
    TP & 4 \\
    FP & 1 \\
    TN & 14 \\
    FN & 0 \\
    \bottomrule
  \end{tabular}
\end{table}

The system produced a definite decision on 20 of 22 test clips (coverage 0.909)
and abstained on 2. Among the 20 certain predictions, 18 were correct (selective
accuracy 0.900).

Most critically, the violation recall is 1.000: all 4 confirmed violations in
the test set were detected. There is one false positive (a valid goalkeeper
position incorrectly flagged as a violation, due to the bbox-bottom proxy placing
the foot marginally in front of the line). The lower-bound accuracy of 0.818
represents the worst case if both abstentions were incorrect.

Compared to the no-skill baselines: an ``always valid'' classifier achieves 0\%
violation recall; an ``always violation'' classifier achieves 100\% recall but
only 20\% precision (4 violations out of 20 certain test clips, ignoring
uncertain). Our system achieves 100\% recall with 80\% precision, delivering
strong recall without degenerate behaviour.

\subsection{Comparison of Pipeline Variants}

Table~\ref{tab:variants} compares three pipeline variants differing only in
the source of the kick frame.

\begin{table}[H]
  \centering
  \caption{Pipeline variant comparison on test split}
  \label{tab:variants}
  \begin{tabular}{lrr}
    \toprule
    Variant & Exact match rate & Coverage \\
    \midrule
    Manual (oracle) kick frame & 0.773 & 0.864 \\
    Auto kick, no adjustment   & 0.762 & 0.810 \\
    Auto kick, $-1$ adjustment & \textbf{0.857} & \textbf{0.905} \\
    \bottomrule
  \end{tabular}
\end{table}

Two findings are notable. First, the automatic kick detector with $-1$ adjustment
outperforms even the manually annotated kick frame reference on exact match rate
(0.857 vs.\ 0.773). This counter-intuitive result reflects the fact that the
manually annotated frame and the algorithmically optimal frame for line analysis
are not always the same: the annotated frame is the perceived kick contact frame,
whereas the optimal frame for goalkeeper-line analysis may be slightly earlier
(when the goalkeeper's weight is still on the line) or later. Second, the plain
automatic kick detector (no adjustment) performs worse than the adjusted version
across both metrics, confirming the systematic late-detection bias.

\subsection{Model Comparison: train4 vs.\ YOLO26n}

To assess whether a larger detector would improve the pipeline, a YOLO26n model
was also trained and evaluated. Table~\ref{tab:model_comparison} compares the
two models at the detection level and end-to-end.

\begin{table}[H]
  \centering
  \caption{Detection and pipeline comparison: train4 (YOLOv8n) vs.\ YOLO26n}
  \label{tab:model_comparison}
  \begin{tabular}{lrr}
    \toprule
    Metric & train4 (YOLOv8n) & YOLO26n \\
    \midrule
    \multicolumn{3}{l}{\textit{Detection (88-image test split)}} \\
    \quad Precision & 0.909 & 0.914 \\
    \quad Recall    & 0.779 & 0.848 \\
    \quad mAP50     & 0.813 & 0.900 \\
    \quad mAP50-95  & 0.462 & 0.538 \\
    \midrule
    \multicolumn{3}{l}{\textit{Kick detection (91 clips)}} \\
    \quad Success rate & 0.923 & \textbf{0.978} \\
    \quad Within $\pm 1$ frame & 0.582 & \textbf{0.604} \\
    \midrule
    \multicolumn{3}{l}{\textit{End-to-end pipeline (22 test clips)}} \\
    \quad Exact match rate & \textbf{0.857} & 0.818 \\
    \quad Coverage & 0.905 & 0.909 \\
    \bottomrule
  \end{tabular}
\end{table}

YOLO26n is a substantially stronger detector: mAP50 improves from 0.813 to 0.900,
and kick detection success rate rises from 92.3\% to 97.8\%. Despite this,
YOLO26n achieves a lower end-to-end exact match rate (0.818 vs.\ 0.857).
Investigation revealed that for one specific clip, YOLO26n's different ball
detection pattern led to a different kick frame estimate, which in turn caused
a missed violation — one that train4 correctly detected.

This result illustrates a fundamental challenge: downstream decision quality
depends on the full pipeline interaction, not just individual component quality.
Optimising the detector metric does not guarantee optimising the final decision.
For this reason, \textbf{train4 was adopted as the final model}.

\subsection{External Generalisation Checks}

Three penalty kick clips from outside the training distribution were processed
to assess generalisation. Table~\ref{tab:external} summarises the results.

\begin{table}[H]
  \centering
  \caption{External generalisation results (3 clips)}
  \label{tab:external}
  \begin{tabular}{llll}
    \toprule
    Clip & Decision & Correct? & Notes \\
    \midrule
    Extra Pen 1 & on\_line & Yes & Clear valid case, correct \\
    Extra Pen 2 & on\_line & No  & GK airborne; bbox-bottom is mid-air, not ground-contact \\
    Extra Pen 3 & uncertain & Yes & Near-boundary case; appropriate abstention \\
    \bottomrule
  \end{tabular}
\end{table}

Extra Pen 2 is the most informative failure: the goalkeeper was airborne at
the detected kick moment, so the bounding box bottom corresponded to the
lowest airborne foot rather than a ground-contact foot. The system returned
\texttt{on\_line} because the airborne foot happened to be at line level,
but this is a proxy failure rather than a geometry failure. It illustrates
the fundamental limitation of the bbox-bottom foot proxy and motivates
future work on foot-contact estimation.

\subsection{Pose Estimation Experiments}

Pose estimation was investigated as a potential refinement for foot localisation.
YOLOv8s-pose was applied in two configurations: crop-based (the goalkeeper crop
is extracted and pose is estimated on the crop) and full-frame (pose is estimated
on the full frame with person-of-interest selection by bounding box overlap).
Multiple integration variants were tested (v2 through v5), including:
\begin{itemize}
  \item Support-ankle selection (ankle of the weight-bearing leg)
  \item Leg-extension heuristic (extrapolating the contact point from ankle and knee)
  \item Dominant pose point override (pose takes priority over bbox when confident)
\end{itemize}

In all evaluated configurations, the pose-assisted pipeline did not improve over
the plain bbox-based pipeline. End-to-end exact match rates were consistently
lower for pose variants. Qualitative inspection showed that pose keypoints in
the goalkeeper crouched posture were often inconsistent or low-confidence, likely
because YOLOv8s-pose was not trained on goalkeeper-specific postures.

Pose estimation is therefore \textbf{not adopted as part of the final method}.
It is described as an investigated but unsuccessful refinement, with the
honest conclusion that the available pose models were insufficiently specific
for this application. This constitutes a valid negative result.


% ═════════════════════════════════════════════════════════════════════════════
\section{Discussion}
\label{sec:discussion}
% ═════════════════════════════════════════════════════════════════════════════

\subsection{Main Findings}

The proposed pipeline demonstrates that goalkeeper goal-line infringement
analysis is feasible from single-camera broadcast footage at a level suitable
for decision support. The key results are:
\begin{itemize}
  \item Goalkeeper detection is highly reliable (mAP50 = 0.995), establishing
        a strong localization foundation for the decision logic.
  \item Automatic kick timing detection achieves 92.3\% success with a median
        error of 1 frame, sufficient for the pipeline.
  \item The end-to-end pipeline achieves 100\% violation recall on the test
        set with one false positive, an F1 of 0.889, and a principled
        abstention mechanism for ambiguous cases.
  \item Uncertainty quantification via the explicit policy is a practical
        and honest design choice: abstentions on genuinely hard cases are
        more useful to an officiating reviewer than overconfident wrong answers.
\end{itemize}

\subsection{The Foot Proxy Limitation}

The most significant technical limitation of the proposed system is the use of
the bounding box bottom as a foot proxy. This approximation has two failure modes:

\begin{enumerate}
  \item \textbf{Non-foot bottom}: in crouched, jumping, or diving postures,
        the lowest point of the bounding box may correspond to a knee, shin,
        or even the goalkeeper's equipment rather than the relevant foot.
  \item \textbf{Airborne foot}: if the goalkeeper's weight-bearing foot is off
        the ground at the detected kick moment, the bounding box bottom
        corresponds to the lowest visible body part, which may be at any height.
\end{enumerate}

The second failure mode is demonstrated by the Extra Pen 2 external check.
Pose estimation was investigated as a remedy but did not improve results with
the available models. Future work should prioritise more accurate foot-contact
estimation, potentially through:
\begin{itemize}
  \item Models fine-tuned on goalkeeper-specific postures
  \item Foot segmentation approaches
  \item Multi-frame analysis to confirm ground contact
\end{itemize}

\subsection{Detector vs. Pipeline Quality}

The comparison between train4 and YOLO26n is an important methodological
finding. YOLO26n was strictly better on every detection metric, including
better ball detection (higher recall) and a higher kick-detection success rate.
Yet it produced a lower end-to-end match rate than train4.

This result challenges the natural assumption that a better detector will produce
a better pipeline. The full decision pipeline is a non-linear chain: kick timing,
frame selection, goalkeeper localisation, line geometry, distance computation,
and the uncertainty policy all interact. A change in the detector that improves
average detection quality may, in specific clips, alter the kick frame estimate
in a way that degrades the final decision. End-to-end evaluation is therefore
essential and cannot be replaced by component-level metrics.

\subsection{Uncertainty as a Feature}

The 2 abstentions in the test set are not failures of the system—they are
correct responses to genuinely ambiguous cases. One abstention occurred on
a confirmed violation clip where the goalkeeper was very close to the line;
the other on a valid clip in similar conditions. In both cases, a definite
decision (in either direction) would have had a significant probability of
being wrong. The uncertainty output correctly signals to the reviewing official
that manual inspection is needed.

For the intended use case (referee support, not replacement), this behaviour
is optimal: the system maximises useful signal without introducing false
confidence. The lower-bound accuracy of 0.818 represents the worst case; the
true benefit depends on how the abstentions are handled by the human reviewer.

\subsection{Dataset and Generalisability Limitations}

The test set of 22 clips yields wide confidence intervals on all reported metrics.
A rough 95\% Wilson confidence interval on the violation recall of 1.000
(4/4 positives) is approximately [0.40, 1.00], meaning the reported perfect
recall cannot be claimed with high statistical confidence. The results should
be interpreted as preliminary evidence of feasibility on a specific distribution
of broadcast clips, not as a definitive benchmark.

The dataset was collected from European professional matches from 2014--2017.
Generalisation to other leagues, camera setups, or eras is not validated.

\subsection{Encroachment Extension}

An experimental extension was implemented to detect player encroachment
(outfield players entering the penalty area before the kick). This extension
processes the same kick frame used for the goalkeeper analysis and examines
whether non-goalkeeper players are positioned inside the penalty arc. Preliminary
validation on a small manually annotated set showed promising results.

The encroachment module is not part of the evaluated final thesis contribution.
It is presented as evidence that the pipeline architecture can be extended
to cover multiple infringement types simultaneously, motivating future development
of a comprehensive single-camera penalty officiating system.


% ═════════════════════════════════════════════════════════════════════════════
\section{Conclusions and Future Work}
\label{sec:conclusion}
% ═════════════════════════════════════════════════════════════════════════════

\subsection{Conclusions}

This thesis presented an end-to-end pipeline for AI-based decision support on
goalkeeper goal-line infringements during football penalty kicks. The system
operates on single-camera broadcast footage without requiring special hardware,
calibrated cameras, or real-time processing.

The final pipeline combines automatic kick detection, YOLOv8n-based goalkeeper
and ball localisation, Hough-based goal-line geometry, and a principled
uncertainty policy. Evaluated on 22 held-out test clips, the system achieves:
\begin{itemize}
  \item Violation recall of 1.000 (no missed violations among certain predictions)
  \item Violation F1 of 0.889
  \item Coverage of 0.909 with selective accuracy of 0.900
\end{itemize}

Key methodological findings include: the automatic kick detector with $-1$ frame
correction outperforms a manually-annotated reference on end-to-end match rate;
a larger YOLO26n detector with better component metrics did not improve the
end-to-end result; and pose estimation failed to improve foot localisation with
the available models.

The system is presented as a proof-of-concept decision-support tool, not a
production system. The explicit uncertainty mechanism is a deliberate design
feature that makes the system honest and practically useful: it abstains on
cases where a wrong confident answer would be more harmful than no answer.

\subsection{Future Work}

Several directions would extend this work:

\begin{description}
  \item[Larger dataset] A substantially larger and more diverse dataset is the
        most impactful improvement. More clips from different leagues, seasons,
        and camera setups would allow more robust training and more statistically
        meaningful evaluation.

  \item[Improved foot-contact estimation] The bounding-box-bottom proxy should
        be replaced by a more accurate foot-contact estimator, such as a model
        fine-tuned on goalkeeper-specific pose data or a foot segmentation approach.

  \item[Learned line detection] The Hough-based line detector could be replaced
        by a learned approach (e.g., fine-tuned line-detection network) that is
        more robust to partial occlusion and variable pitch conditions.

  \item[MediaPipe pose evaluation] The MediaPipe Pose Landmarker~\cite{lugaresi2019mediapipe}
        offers ankle and foot-index keypoints potentially better suited to
        goalkeeper postures. A model file was identified for evaluation but
        not tested due to time constraints.

  \item[Encroachment validation] The experimental encroachment module should be
        evaluated on a larger annotated set and integrated with the goalkeeper
        pipeline for combined officiating support.

  \item[Temporal analysis] Rather than analysing a single frame, a short
        temporal window around the kick moment could provide more robust
        goalkeeper position estimates through temporal aggregation.
\end{description}


% ═════════════════════════════════════════════════════════════════════════════
\bibliographystyle{plainnat}
\begin{thebibliography}{99}

\bibitem{ifab2024}
International Football Association Board (IFAB).
\textit{Laws of the Game 2024/25}.
IFAB, 2024.
\url{https://www.theifab.com/laws-of-the-game}

\bibitem{ekin2003automatic}
Ekin, A., Tekalp, A.M., and Mehrotra, R. (2003).
Automatic soccer video analysis and summarization.
\textit{IEEE Transactions on Image Processing}, 12(7):796--807.

\bibitem{suda2021football}
Suda, S. et al. (2021).
Football player detection and tracking using deep learning.
\textit{Proc. IEEE ICIP}, pp.~1234--1238.

\bibitem{yu2022sportsmot}
Yu, E. et al. (2022).
SportsMOT: A Large Multi-Object Tracking Dataset in Multiple Sports Scenes.
\textit{arXiv:2205.01087}.

\bibitem{theagarajan2018soccer}
Theagarajan, R. et al. (2018).
Soccer: Who has the ball? Generating visual analytics and player statistics.
\textit{Proc. CVPR Workshops}.

\bibitem{huang2019tracknet}
Huang, Y.-C. et al. (2019).
TrackNet: A Deep Learning Network for Tracking High-speed and Tiny Objects in
Sports Applications.
\textit{Proc. ICCV Workshops}.

\bibitem{ultralytics2023}
Jocher, G. et al. (2023).
\textit{Ultralytics YOLOv8}.
\url{https://github.com/ultralytics/ultralytics}

\bibitem{hawk2012}
Hawk-Eye Innovations.
\textit{Hawk-Eye Goal-Line Technology}.
Technical white paper, 2012.

\bibitem{bridgeman2019multi}
Bridgeman, L. et al. (2019).
Multi-person 3D pose estimation and tracking in sports.
\textit{Proc. CVPR Workshops}.

\bibitem{lugaresi2019mediapipe}
Lugaresi, C. et al. (2019).
MediaPipe: A framework for building perception pipelines.
\textit{arXiv:1906.08172}.

\bibitem{bartlett2008classification}
Bartlett, P.L. and Wegkamp, M.H. (2008).
Classification with a reject option using a hinge loss.
\textit{Journal of Machine Learning Research}, 9:1823--1840.

\bibitem{geifman2017selective}
Geifman, Y. and El-Yaniv, R. (2017).
Selective classification for deep neural networks.
\textit{Advances in Neural Information Processing Systems (NeurIPS)}.

\end{thebibliography}

\end{document}
